% Figure study_area: generated by figures/rse_fig_study_area.py; placed where the text first discusses it.
\begin{figure}[!tbp]
  \centering
  \includegraphics[width=\textwidth]{../figures/out/rse_fig_study_area.pdf}
  \caption{\textbf{Study area.}
  (a)~The North Caspian study domain (thick black outline) and its seven reporting sectors (thin outlines), each named on the map.
  A sector counts as checked (bold name in a black box) when, in its ICESat-2 test blocks, the mean and the median of the waterline elevation model minus the ICESat-2 ATL08 heights, and the same two statistics against ATL06, all lie within $\pm\metric{params.checked_tol_m}$~m.
  The Volga delta and Komsomolets Bay fail this rule and the Ural delta has no test block; these three sectors are hatched and their names are set in italics.
  Vermilion squares: the \metric{h2.atl08.n_test_blocks} test blocks of \metric{h2.block_size_km}~km that hold the ICESat-2 points used to test the elevation model.
  Tan: seabed that emerged between the \metric{params.highstand_label} high stand and the \metric{params.lowstand_label} low stand (water in the Landsat extent of the first period and not in that of the second).
  Light blue: the open sea of the \metric{params.lowstand_label} low stand, from Landsat.
  Dash-dotted line: international boundary (Natural Earth).
  Lambert azimuthal equal-area projection centred at \metric{params.laea_lat0_deg}\textdegree{}N, \metric{params.laea_lon0_deg}\textdegree{}E; graticule every \metric{params.fig_graticule_deg}\textdegree{}; scale bar \metric{params.fig_scale_bar_km}~km.
  (b)~The study domain (vermilion outline) in the Caspian region (Natural Earth land, lakes and boundaries).}
  \label{fig:rse_study_area}
\end{figure}
